\documentclass[11pt,a4paper,sans]{moderncv}

% Style options: 'casual', 'classic', 'banking', 'oldstyle', 'fancy'
\moderncvstyle{classic} 
\font\myfont=cmr12 at 10pt
% Color options: 'blue', 'orange', 'green', 'red', 'purple', 'grey', 'black'
\moderncvcolor{blue} 
\renewcommand{\familydefault}{helvetica}
\usepackage[utf8]{inputenc}

% Adjust the page margins
\usepackage[scale=0.75, top=2cm, bottom = 2cm]{geometry}
\usepackage[ngerman]{babel} 
% Personal Data
\name{Samuel}{O'Neill}
\title{{\myfont Informatikstudent an der ETH Zürich}}                               
\address{Haldenpark 22}{4313, Möhlin}{Switzerland}
\phone[mobile]{+41 79 868 58 84}
\email{samueloneill.dev@gmail.com}
\social[linkedin][www.linkedin.com/in/samuel-o-neill-98a1b01b8]{Samuel O'Neill}
\social[github]{Diddleydope}
\photo[110px][2px]{samuel_picture.png}

%----------------------------------------------------------------------------------
%            Content
%----------------------------------------------------------------------------------
\begin{document}
\makecvtitle

\section{Ausbildung}
\cventry{2019-2023}{Gymnasium Muttenz}{Schwerpunkt Mathe und Physik}{Muttenz}{}{}
\cventry{2024-2027}{Bachelor der Informatik}{ETH Zürich}{Zürich}{Notenschnitt: 4.9 (Drittes Sem.)}{Besuche zurzeit  Kernfächer wie Introduction to Machine Learning und Visual Computing}

\section{Erfahrung und Projekte}
\cventry{2022-2023}{Dezentralisierter Messenger}{Maturaarbeit}{Note: 5,5}{\newline Sprachen/Technologien: Javascript, CSS, HTML, GUN.js}{
Ein dezentralisierter Messenger mit GUN.js, welcher Verschlüsselung nutzt, um durch Dezentralisierung einen sicheren Nachrichtenaustausch zu ermöglichen.}
\cventry{2024 (Jan - Mai)}{Praktikum}{SPECS Zürich GmbH}{\newline Sprachen/Technologien: Python, PyPi}{}{Entwicklung von zwei Python-Paketen zur Kommunikation mit Nanonis Tramea und SPM-Controllern via Python (Siehe nanonis-tramea (35x im letzten Monat heruntergeladen) und nanonis-spm (174x im letzten Monat heruntergeladen) auf pypi)
}

\cventry{2024-Today}{ShinyDiary | Persönliches Projekt}{\newline Sprachen/Technologien: Rust, Tauri, Svelte, TypeScript, Firebase}{}{}{Entwicklung einer plattformübergreifenden Desktop-Anwendung für Pokémon „Shiny Hunting“-Enthusiasten zur Nachverfolgung, Protokollierung und Automatisierung von Begegnungsstatistiken. Implementierung einer „debounced“ Datenpersistenz-Strategie zur verbesserten Skalierbarkeit und Kostensenkung. (Siehe GitHub)}
\section{Languages}
\cvitemwithcomment{Englisch}{Fliessend}{Muttersprache}
\cvitemwithcomment{Deutsch}{Fliessend}{Muttersprache}
\cvitemwithcomment{Französisch}{Fortgeschritten}{B2 Level}

\section*{Technische Fähigkeiten}
\noindent \textbf{Sprachen:} Python, C, C++, Java, Haskell, JavaScript, SQL, HTML/CSS \\
\textbf{Machine Learning:} PyTorch, Scikit-learn, Computer Vision \\
\textbf{Web Development:} Svelte, Node.js, Tauri  \\
\textbf{Developer Tools:} Git, Docker, Kathara, VS Code, Linux, GitHub Actions



\end{document}